% Options for packages loaded elsewhere
\PassOptionsToPackage{unicode}{hyperref}
\PassOptionsToPackage{hyphens}{url}
\PassOptionsToPackage{dvipsnames,svgnames,x11names}{xcolor}
\documentclass[
  11pt,
  a4paper,
]{article}
\usepackage{xcolor}
\usepackage[margin=1in]{geometry}
\usepackage{amsmath,amssymb}
\setcounter{secnumdepth}{-\maxdimen} % remove section numbering
\usepackage{iftex}
\ifPDFTeX
  \usepackage[T1]{fontenc}
  \usepackage[utf8]{inputenc}
  \usepackage{textcomp} % provide euro and other symbols
\else % if luatex or xetex
  \usepackage{unicode-math} % this also loads fontspec
  \defaultfontfeatures{Scale=MatchLowercase}
  \defaultfontfeatures[\rmfamily]{Ligatures=TeX,Scale=1}
\fi
\usepackage{lmodern}
\ifPDFTeX\else
  % xetex/luatex font selection
  \ifLuaTeX
    \usepackage{luaotfload}
    \directlua{luaotfload.add_fallback("mainfontfallback",{
      "NotoSerifBengali:mode=harf"
    })}
  \fi
  \setmainfont[,RawFeature={fallback=mainfontfallback}]{Latin Modern Roman}
  \ifLuaTeX
    \usepackage{luaotfload}
    \directlua{luaotfload.add_fallback("sansfontfallback",{
      "NotoSansBengali:mode=harf"
    })}
  \fi
  \setsansfont[,RawFeature={fallback=sansfontfallback}]{Latin Modern Sans}
  \ifLuaTeX
    \usepackage{luaotfload}
    \directlua{luaotfload.add_fallback("monofontfallback",{
      "NotoSansBengali:mode=harf"
    })}
  \fi
  \setmonofont[,RawFeature={fallback=monofontfallback}]{DejaVu Sans Mono}
\fi
% Use upquote if available, for straight quotes in verbatim environments
\IfFileExists{upquote.sty}{\usepackage{upquote}}{}
\IfFileExists{microtype.sty}{% use microtype if available
  \usepackage[]{microtype}
  \UseMicrotypeSet[protrusion]{basicmath} % disable protrusion for tt fonts
}{}
\makeatletter
\@ifundefined{KOMAClassName}{% if non-KOMA class
  \IfFileExists{parskip.sty}{%
    \usepackage{parskip}
  }{% else
    \setlength{\parindent}{0pt}
    \setlength{\parskip}{6pt plus 2pt minus 1pt}}
}{% if KOMA class
  \KOMAoptions{parskip=half}}
\makeatother
\usepackage{longtable,booktabs,array}
\usepackage{calc} % for calculating minipage widths
% Correct order of tables after \paragraph or \subparagraph
\usepackage{etoolbox}
\makeatletter
\patchcmd\longtable{\par}{\if@noskipsec\mbox{}\fi\par}{}{}
\makeatother
% Allow footnotes in longtable head/foot
\IfFileExists{footnotehyper.sty}{\usepackage{footnotehyper}}{\usepackage{footnote}}
\makesavenoteenv{longtable}
\setlength{\emergencystretch}{3em} % prevent overfull lines
\providecommand{\tightlist}{%
  \setlength{\itemsep}{0pt}\setlength{\parskip}{0pt}}
% paper_preamble.tex — shared LaTeX preamble for the PCMA/SI paper set.
%
% Compiled with lualatex (NOT pdflatex) because several papers contain
% Bengali script: the CLL paper realises one cognitive state into
% English, Bengali and Banglish, so the Bengali output is EVIDENCE and
% cannot be transliterated away without destroying the point being made.
%
% Font selection is NOT done here. pandoc >= 3 has native
% `mainfontfallback` / `monofontfallback` variables that emit the
% luaotfload fallback registration in the correct position relative to
% unicode-math; doing it by hand in this header lands after unicode-math
% and breaks font loading outright. See tools/build_papers.sh.

\usepackage{amsmath}
\usepackage{amssymb}
\usepackage{mathtools}   % \xleftrightarrow etc. used in cll.md §5
\usepackage{booktabs}

% ---------------------------------------------------------------------
% Unicode symbols appearing in PROSE rather than inside $...$.
%
% The papers were written as markdown and read as markdown, so they use
% real characters — "the ω-boost", "α-child", "ΔR coherence" — instead of
% LaTeX macros. Under unicode-math those codepoints are math-mode-only
% and abort the build ("\symbballowed only in math mode"). Mapping each
% to \ensuremath keeps the markdown readable on GitHub AND compilable,
% which is the whole point of keeping one source.
% ---------------------------------------------------------------------
\usepackage{newunicodechar}

% Greek, lower case
\newunicodechar{α}{\ensuremath{\alpha}}
\newunicodechar{β}{\ensuremath{\beta}}
\newunicodechar{γ}{\ensuremath{\gamma}}
\newunicodechar{δ}{\ensuremath{\delta}}
\newunicodechar{ε}{\ensuremath{\varepsilon}}
\newunicodechar{η}{\ensuremath{\eta}}
\newunicodechar{θ}{\ensuremath{\theta}}
\newunicodechar{κ}{\ensuremath{\kappa}}
\newunicodechar{λ}{\ensuremath{\lambda}}
\newunicodechar{μ}{\ensuremath{\mu}}
\newunicodechar{π}{\ensuremath{\pi}}
\newunicodechar{ρ}{\ensuremath{\rho}}
\newunicodechar{σ}{\ensuremath{\sigma}}
\newunicodechar{τ}{\ensuremath{\tau}}
\newunicodechar{φ}{\ensuremath{\varphi}}
\newunicodechar{χ}{\ensuremath{\chi}}
\newunicodechar{ψ}{\ensuremath{\psi}}
\newunicodechar{ω}{\ensuremath{\omega}}
\newunicodechar{ν}{\ensuremath{\nu}}
\newunicodechar{ξ}{\ensuremath{\xi}}
\newunicodechar{ζ}{\ensuremath{\zeta}}
\newunicodechar{ι}{\ensuremath{\iota}}

% Greek, upper case
\newunicodechar{Δ}{\ensuremath{\Delta}}
\newunicodechar{Φ}{\ensuremath{\Phi}}
\newunicodechar{Σ}{\ensuremath{\Sigma}}
\newunicodechar{Ψ}{\ensuremath{\Psi}}
\newunicodechar{Ω}{\ensuremath{\Omega}}
\newunicodechar{Γ}{\ensuremath{\Gamma}}
\newunicodechar{Λ}{\ensuremath{\Lambda}}
\newunicodechar{Θ}{\ensuremath{\Theta}}
\newunicodechar{Π}{\ensuremath{\Pi}}

% Operators and relations
\newunicodechar{×}{\ensuremath{\times}}
\newunicodechar{·}{\ensuremath{\cdot}}
\newunicodechar{±}{\ensuremath{\pm}}
\newunicodechar{−}{\ensuremath{-}}
\newunicodechar{≈}{\ensuremath{\approx}}
\newunicodechar{≤}{\ensuremath{\leq}}
\newunicodechar{≥}{\ensuremath{\geq}}
\newunicodechar{≠}{\ensuremath{\neq}}
\newunicodechar{∈}{\ensuremath{\in}}
\newunicodechar{∂}{\ensuremath{\partial}}
\newunicodechar{∇}{\ensuremath{\nabla}}
\newunicodechar{√}{\ensuremath{\sqrt{}}}
\newunicodechar{∞}{\ensuremath{\infty}}
\newunicodechar{∝}{\ensuremath{\propto}}
\newunicodechar{∑}{\ensuremath{\sum}}
\newunicodechar{∫}{\ensuremath{\int}}
\newunicodechar{ℝ}{\ensuremath{\mathbb{R}}}
\newunicodechar{ℕ}{\ensuremath{\mathbb{N}}}
\newunicodechar{‖}{\ensuremath{\|}}
\newunicodechar{⊕}{\ensuremath{\oplus}}
\newunicodechar{⊗}{\ensuremath{\otimes}}
\newunicodechar{∀}{\ensuremath{\forall}}
\newunicodechar{∃}{\ensuremath{\exists}}
\newunicodechar{¬}{\ensuremath{\neg}}
\newunicodechar{∧}{\ensuremath{\wedge}}
\newunicodechar{∨}{\ensuremath{\vee}}

% Arrows
\newunicodechar{→}{\ensuremath{\rightarrow}}
\newunicodechar{←}{\ensuremath{\leftarrow}}
\newunicodechar{↔}{\ensuremath{\leftrightarrow}}
\newunicodechar{↑}{\ensuremath{\uparrow}}
\newunicodechar{↓}{\ensuremath{\downarrow}}
\newunicodechar{⇒}{\ensuremath{\Rightarrow}}
\newunicodechar{⇐}{\ensuremath{\Leftarrow}}
\newunicodechar{⇔}{\ensuremath{\Leftrightarrow}}

% Super/subscripts written as literal characters
\newunicodechar{²}{\ensuremath{^{2}}}
\newunicodechar{³}{\ensuremath{^{3}}}
\newunicodechar{⁴}{\ensuremath{^{4}}}
\newunicodechar{⁵}{\ensuremath{^{5}}}
\newunicodechar{⁶}{\ensuremath{^{6}}}
\newunicodechar{₀}{\ensuremath{_{0}}}
\newunicodechar{₁}{\ensuremath{_{1}}}
\newunicodechar{₂}{\ensuremath{_{2}}}
\newunicodechar{₃}{\ensuremath{_{3}}}
\newunicodechar{ᵀ}{\ensuremath{^{\mathsf{T}}}}

% Status glyphs. The papers use emoji as status markers in tables
% (implemented / failed / warning). Mapped to typographic equivalents:
% a colour emoji font in lualatex is unreliable, and the MEANING is what
% carries — a reader needs to see "this one failed", not a red cross.
\newunicodechar{✅}{\ensuremath{\checkmark}}
\newunicodechar{✓}{\ensuremath{\checkmark}}
\newunicodechar{❌}{\ensuremath{\times}}
\newunicodechar{✗}{\ensuremath{\times}}
\newunicodechar{⚠}{\textbf{!}}
\newunicodechar{🧩}{\textbf{[piece]}}
\newunicodechar{📋}{\textbf{[list]}}
\newunicodechar{🟡}{\textbf{[partial]}}
\newunicodechar{🔴}{\textbf{[dead]}}
\newunicodechar{🟢}{\textbf{[live]}}
\newunicodechar{▼}{\ensuremath{\blacktriangledown}}
\newunicodechar{►}{\ensuremath{\blacktriangleright}}
\newunicodechar{□}{\ensuremath{\square}}
\newunicodechar{█}{\rule{0.6em}{0.9ex}}
\newunicodechar{•}{\textbullet}
\newunicodechar{†}{\textdagger}
\newunicodechar{‑}{-}
\newunicodechar{‐}{-}
\newunicodechar{…}{\ldots}
\newunicodechar{।}{\textbar}   % danda: rendered by the Bengali fallback in text
\newunicodechar{⁻}{\ensuremath{^{-}}}

% Zero-width formatting codepoints. These are invisible in an editor but
% are real characters to LaTeX: an emoji written as "⚠️" is TWO
% codepoints (U+26A0 U+FE0F), so declaring only the visible one leaves a
% variation selector with no definition and the build dies on a symbol
% nobody can see. Both are mapped to nothing.
\newunicodechar{^^^^fe0f}{}    % VARIATION SELECTOR-16
\newunicodechar{^^^^200d}{}    % ZERO WIDTH JOINER
\newunicodechar{^^^^200b}{}    % ZERO WIDTH SPACE

\usepackage{longtable}
\usepackage{array}

% Comparison matrices and cost ledgers are wide. Let them shrink to the
% text block rather than overflow the margin.
\usepackage{adjustbox}

% Code blocks: wrap long lines instead of running off the page. Several
% papers quote real C++ with long identifiers.
\usepackage{fvextra}
\fvset{breaklines=true,breakanywhere=true,fontsize=\small}

\setlength{\parskip}{0.5em}
\setlength{\parindent}{0pt}

\usepackage{fancyhdr}
\pagestyle{fancy}
\fancyhf{}
\fancyfoot[C]{\thepage}
\fancyhead[L]{\small\itshape Synthetic Intelligence — research lineage}
\fancyhead[R]{\small\itshape irfan@furylogic.com}
\renewcommand{\headrulewidth}{0.2pt}
\setlength{\headheight}{14pt}
\setlength{\headsep}{18pt}
\usepackage{bookmark}
\IfFileExists{xurl.sty}{\usepackage{xurl}}{} % add URL line breaks if available
\urlstyle{same}
\hypersetup{
  pdftitle={The Cognitive Multimodal Engine (CME): A Unified Non-Embedding Architecture for Real-Time Multi-Modal Understanding via Physics-Based Cognitive State Evolution},
  pdfauthor={Irfan Mahir --- Independent Researcher, Dhaka, Bangladesh --- irfan@furylogic.com},
  colorlinks=true,
  linkcolor={black},
  filecolor={Maroon},
  citecolor={Blue},
  urlcolor={[HTML]{0B5394}},
  pdfcreator={LaTeX via pandoc}}

\title{The Cognitive Multimodal Engine (CME): A Unified Non-Embedding
Architecture for Real-Time Multi-Modal Understanding via Physics-Based
Cognitive State Evolution}
\author{Irfan Mahir --- Independent Researcher, Dhaka, Bangladesh ---
irfan@furylogic.com}
\date{}

\begin{document}
\maketitle

{
\hypersetup{linkcolor=}
\setcounter{tocdepth}{3}
\tableofcontents
}
\section{The Cognitive Multimodal Engine (CME): A Unified Non-Embedding
Architecture for Real-Time Multi-Modal Understanding via Physics-Based
Cognitive State
Evolution}\label{the-cognitive-multimodal-engine-cme-a-unified-non-embedding-architecture-for-real-time-multi-modal-understanding-via-physics-based-cognitive-state-evolution}

\textbf{Irfan Mahir}\\
\emph{Independent Researcher, Dhaka, Bangladesh}\\
irfan@furylogic.com

\begin{quote}
\emph{``Thinking and understanding is the art of innovation and
invention.''}
\end{quote}

\begin{center}\rule{0.5\linewidth}{0.5pt}\end{center}

\subsection{Abstract}\label{abstract}

We present the \textbf{Cognitive Multimodal Engine (CME)} , a unified
architecture for real-time multi-modal understanding that fundamentally
departs from both the token-prediction paradigm of Transformers and the
embedding-based fusion of conventional multimodal systems. Unlike
systems such as Gemini Live that process each modality through separate
encoders and fuse them via cross-attention over token sequences, CME
maps all modalities---text, audio, image, video, and live
streams---directly into a single fixed-size cognitive state
\(c \in \mathbb{R}^{d}\) that evolves continuously according to
Hamiltonian dynamics.

CME integrates three core innovations: (1) a \textbf{modality-agnostic
byte-level encoder} that maps raw input directly to the cognitive state
without tokenization or embeddings; (2) a \textbf{Semantic Nostalgia
Buffer} that maintains compact hash-based fingerprints for O(1)
cross-temporal pattern matching across all modalities; and (3) a
\textbf{Cognitive State Manifold} governed by Hamiltonian dynamics,
which replaces the autoregressive token stream with continuous
physics-based reasoning over a fixed-size cognitive state.

For streaming video, CME incorporates a \textbf{Differential Frame
Significance (DFS)} pre-processor that operates at the pixel-vector flow
level to detect semantically meaningful changes before any encoding
occurs, achieving a 78.3\% latency reduction over state-of-the-art
online VideoLLMs. For text, audio, and images, the same cognitive state
evolution framework applies without modification---the only difference
is the input encoder.

Crucially, CME operates without tokenization, without embeddings,
without KV-cache, and without any form of sequence memory---replacing
them with a single O(1) cognitive state that compresses the entire
interaction history. The architecture is grounded in the Cognitive
Continuum Hypothesis (CCH), which establishes reasoning as geodesic flow
on a logical manifold with a conservative vector field
\(F = -\nabla\Phi\) guaranteeing \(\nabla \times F = 0\) and thus 100\%
logical transitivity.

Our work establishes a new paradigm for multi-modal AI that moves beyond
``separate encoders with cross-attention'' toward truly unified,
continuous, event-driven intelligence---all without relying on
embeddings, tokenizers, or autoregressive generation.

\begin{center}\rule{0.5\linewidth}{0.5pt}\end{center}

\subsection{1. Introduction}\label{introduction}

The proliferation of multi-modal AI systems---from AI assistants and
autonomous systems to real-time translation and live streaming
services---has created an urgent need for models capable of continuous,
low-latency understanding across text, audio, image, video, and live
streams. However, the dominant paradigm of multi-modal AI, exemplified
by systems like Gemini Live, remains fundamentally grounded in the
Transformer architecture: each modality is processed through a separate
encoder, mapped to token sequences, and fused via cross-attention.

This paradigm carries three irreducible costs:

\begin{enumerate}
\def\labelenumi{\arabic{enumi}.}
\tightlist
\item
  \textbf{Separate Embedding Spaces:} Each modality is mapped to its own
  embedding space, requiring complex fusion mechanisms and preventing
  true cross-modal understanding.
\item
  \textbf{Tokenization Overhead:} Every modality must be tokenized,
  introducing information loss and arbitrary segmentation.
\item
  \textbf{KV-Cache Memory:} The growing sequence of tokens imposes O(N)
  memory complexity, eventually leading to out-of-memory errors or
  catastrophic forgetting.
\end{enumerate}

Recent work on multi-modal LLMs such as LLaVA-NeXT, VideoLLaMA, and
Gemini has made significant progress. However, all such approaches treat
multi-modality as a problem of \emph{fusion} rather than
\emph{unification}. They assume that each modality must be processed
separately and then combined, rather than recognizing that understanding
itself is modality-agnostic.

We argue that the unit of processing for multi-modal AI should not be
the token, the frame, or the embedding---it should be the
\emph{cognitive state}. The human brain does not maintain separate
embedding spaces for vision, language, and sound. It maintains a unified
cognitive state that integrates information across modalities and
evolves continuously through time.

This insight is the foundation of our work.

We introduce the \textbf{Cognitive Multimodal Engine (CME)} , which
operationalizes this principle through three synergistic innovations:

\begin{enumerate}
\def\labelenumi{\arabic{enumi}.}
\item
  \textbf{Modality-Agnostic Byte-Level Encoding:} All inputs---text,
  audio, images, video, live streams---are mapped directly to a
  fixed-size cognitive state \(c\) via a lightweight byte-level encoder.
  No tokenization, no embeddings, no modality-specific feature
  extraction.
\item
  \textbf{Semantic Nostalgia Buffer:} Maintains compact, hash-based
  semantic fingerprints of historical context, enabling O(1) lookup for
  whether a given semantic content has been seen before. This replaces
  the growing KV-cache with a fixed-size associative memory, analogous
  to human episodic memory.
\item
  \textbf{Cognitive State Manifold with Hamiltonian Dynamics:} Instead
  of processing inputs as sequences of tokens, CME maintains a
  fixed-size cognitive state \(c \in \mathbb{R}^{d}\) that evolves
  continuously according to Hamiltonian dynamics. This state compresses
  the entire interaction history into O(1) memory and serves as the sole
  representation for reasoning, which is formulated as energy
  minimization over a logical manifold.
\end{enumerate}

CME is built upon the \textbf{Predictive Cognitive Manifold Architecture
(P-CMA)} and is grounded in the \textbf{Cognitive Continuum Hypothesis
(CCH)} , which establishes reasoning as geodesic flow on a logical
manifold with a conservative vector field \(F = -\nabla\Phi\)
guaranteeing \(\nabla \times F = 0\) and thus 100\% logical
transitivity. Unlike Transformers, CME has no tokens, no KV-cache, no
context window, no autoregressive generation, and no embeddings---only a
continuously evolving cognitive state.

\begin{center}\rule{0.5\linewidth}{0.5pt}\end{center}

\subsection{2. Related Work}\label{related-work}

\subsubsection{2.1 Multi-Modal LLMs}\label{multi-modal-llms}

Recent advances in multi-modal LLMs have primarily focused on extending
Transformer-based architectures to handle multiple modalities.
\textbf{LLaVA-NeXT} uses a visual encoder and a language model with
projection layers. \textbf{VideoLLaMA} extends this to video with
frame-level encoders and temporal modeling. \textbf{Qwen-VL} uses a
vision transformer with a language model. \textbf{Gemini} uses separate
encoders for text, image, audio, and video, with cross-attention for
fusion. \textbf{GPT-4 with Vision} uses a vision encoder with a language
model.

These approaches share a common limitation: they rely on separate
embedding spaces for each modality, requiring complex cross-attention
mechanisms and limiting true cross-modal understanding. They also
inherit the Transformer's O(N) memory complexity and autoregressive
generation constraints.

\subsubsection{2.2 Online Streaming Video
Understanding}\label{online-streaming-video-understanding}

Recent work on online video understanding includes
\textbf{VideoLLM-online}, which proposes the LIVE framework for
real-time conversation. \textbf{TimeChat-Online} introduces Differential
Token Drop (DTD), which reduces tokens by 82.8\% while maintaining 98\%
performance. \textbf{StreamMind} introduces event-gated LLM invocation.
\textbf{StreamForest} introduces Persistent Event Memory Forest (PEMF).
\textbf{StreamBridge} transforms offline Video-LLMs into
streaming-capable models. \textbf{LiveStar} addresses real-time
responsiveness and narrative coherence. \textbf{LiveVLM} employs
streaming-oriented KV cache and retrieval.

All these methods remain within the frame-based, token-predicting
paradigm. None has integrated event detection before encoding, nor
replaced the token stream with a physics-based cognitive state.

\subsubsection{2.3 Event-Gated and Memory-Augmented
Approaches}\label{event-gated-and-memory-augmented-approaches}

Recent work on event-gated processing includes \textbf{EventSTU}, which
proposes an event-guided framework. \textbf{Video-EM} presents an
event-centric episodic memory framework. \textbf{EventGPT} integrates
LLMs with event-based vision. \textbf{LLaFEA} uses frame-event
complementary fusion. \textbf{GCAgent} employs schematic and narrative
episodic memory. \textbf{HEM-LLM} uses hierarchical event-based memory.
\textbf{CausalMem} constructs a dynamic visual memory update mechanism.
\textbf{MA-LMM} uses a memory bank for long-term video understanding.
\textbf{HERMES} introduces episodes and semantics for long-form
understanding. \textbf{CurveStream} uses curvature-aware hierarchical
visual memory management. \textbf{WAT} builds a hierarchical memory
system with short-term and long-term memory.

These approaches demonstrate the growing recognition that memory and
event detection are critical for streaming understanding. However, they
remain additive components attached to frame-based, token-predicting
architectures.

\subsubsection{2.4 Non-Transformer Multimodal
Architectures}\label{non-transformer-multimodal-architectures}

Recent work exploring non-Transformer backbones includes
\textbf{ModRWKV}, which builds upon the RWKV7 architecture for
multi-source information fusion. \textbf{NEO} claims to be the world's
first native multimodal architecture. \textbf{ChronoMem} enables agents
to store, consolidate, and retrieve multimodal experiences.
\textbf{Lyra} provides an efficient speech-centric framework for
omni-cognition.

While these approaches demonstrate that non-Transformer architectures
are viable for multimodal tasks, they still rely on tokenization,
embeddings, and sequence-based processing. CME goes further by
eliminating tokenization, embeddings, and sequence processing entirely.

\begin{center}\rule{0.5\linewidth}{0.5pt}\end{center}

\subsection{3. Methodology}\label{methodology}

\subsubsection{3.1 Problem Formulation}\label{problem-formulation}

Let the system receive a sequence of multi-modal inputs
\(\{x_1, x_2, \dots, x_t, \dots\}\) where each \(x_t\) can be text,
audio, image, video frame, or a combination. A multi-modal understanding
system must process this unbounded stream and generate responses \(y_t\)
without access to future inputs.

Traditional approaches process each input through modality-specific
encoders, producing embeddings that are then tokenized and processed by
an autoregressive LLM.

CME departs from this formulation entirely. Instead of processing inputs
as sequences of tokens, CME maintains a fixed-size cognitive state
\(c_t \in \mathbb{R}^{d}\) that evolves according to:

\[
c_{t+1} = \mathcal{H}\left(c_t, \Delta(x_t, x_{t-1})\right)
\]

where \(\Delta(x_t, x_{t-1})\) is a measure of semantic change between
inputs, and \(\mathcal{H}\) is the Hamiltonian dynamics operator.

\subsubsection{3.2 Modality-Agnostic Byte-Level
Encoding}\label{modality-agnostic-byte-level-encoding}

All inputs are encoded directly from raw bytes to the cognitive state
without tokenization or embeddings. This is the key innovation that
distinguishes CME from all prior work: there is no modality-specific
feature extraction, no tokenization, no BPE, no embeddings.

For each modality, the encoding is:

\textbf{Text:} Raw UTF-8 bytes → MLP encoder →
\(z_t \in \mathbb{R}^{d}\)

\textbf{Audio:} Raw PCM bytes → 1D CNN → \(z_t \in \mathbb{R}^{d}\)

\textbf{Image:} Raw pixel bytes → lightweight CNN →
\(z_t \in \mathbb{R}^{d}\)

\textbf{Video:} Raw frame bytes → DFS pre-processor + Nostalgia Buffer →
\(z_t \in \mathbb{R}^{d}\)

\textbf{Live Stream:} Any of the above → continuous updates to
\(z_t \in \mathbb{R}^{d}\)

In all cases, the output is the same: a continuous vector
\(z_t \in \mathbb{R}^{d}\) that represents the semantic content of the
input, mapped directly to the cognitive manifold without any
tokenization or embedding step.

\subsubsection{3.3 Differential Frame Significance (DFS)
Pre-processor}\label{differential-frame-significance-dfs-pre-processor}

For video inputs, the DFS pre-processor computes frame-to-frame
differences to identify semantically meaningful changes \emph{before}
any encoding occurs.

\textbf{Pixel-Vector Flow Computation:} For each incoming frame \(f_t\),
we compute the optical flow magnitude
\(m_t = \|\text{Flow}(f_t, f_{t-1})\|\) and histogram difference
\(h_t = \text{HistDiff}(f_t, f_{t-1})\). These are combined into a
significance score:

\[
s_t = \alpha \cdot \bar{m}_t + (1 - \alpha) \cdot \bar{h}_t
\]

where \(\bar{m}_t\) and \(\bar{h}_t\) are normalized values, and
\(\alpha\) is a weighting parameter.

\textbf{Adaptive Thresholding:} Rather than a fixed threshold, DFS
employs an adaptive threshold \(\tau_t = \mu_s + \beta \cdot \sigma_s\),
where \(\mu_s\) and \(\sigma_s\) are the running mean and standard
deviation of significance scores.

\textbf{Graded Significance:} Unlike binary event detection, DFS
produces a continuous score \(s_t \in [0, 1]\) that determines the
degree of cognitive engagement:

\[
g_t = \begin{cases}
0 & s_t < \tau_{\text{low}} \\
\frac{s_t - \tau_{\text{low}}}{\tau_{\text{high}} - \tau_{\text{low}}} & \tau_{\text{low}} \leq s_t \leq \tau_{\text{high}} \\
1 & s_t > \tau_{\text{high}}
\end{cases}
\]

Frames with \(g_t = 0\) are skipped entirely---no encoding, no
tokenization, no LLM invocation.

\subsubsection{3.4 Semantic Nostalgia
Buffer}\label{semantic-nostalgia-buffer}

The Nostalgia Buffer maintains compact, hash-based semantic fingerprints
for rapid cross-temporal matching across all modalities.

\textbf{Semantic Fingerprint Generation:} For any input with
significance \(s_t > \tau_{\text{low}}\), we generate a fingerprint
\(p_t = \text{Hash}(E_{\text{light}}(x_t))\), where \(E_{\text{light}}\)
is a lightweight feature extractor (e.g., a small CNN with \textless1M
parameters). The fingerprint is a compact representation that captures
the semantic essence of the input without storing the full features.

\textbf{Temporal Hashing:} The buffer maintains a hash map
\(\mathcal{H}: \text{Hash} \to \{(t, s_t, \text{metadata})\}\), enabling
\(O(1)\) lookup for whether we have seen this semantic content before.
When a new input arrives, the system queries:

\[
\text{exists} = \mathcal{H}.\text{contains}(p_t)
\]

\textbf{Nostalgia Retrieval:} If the semantic content has been seen
before, the system retrieves the corresponding cognitive state
contribution from the buffer, avoiding redundant processing. This is
analogous to human episodic memory, where repeated experiences are
recognized without re-processing.

\subsubsection{3.5 Cognitive State Manifold and Hamiltonian
Dynamics}\label{cognitive-state-manifold-and-hamiltonian-dynamics}

CME is built upon the \textbf{Predictive Cognitive Manifold Architecture
(P-CMA)} and the \textbf{Cognitive Continuum Hypothesis (CCH)} . The
cognitive state \(c \in \mathbb{R}^{d}\) evolves according to
Hamiltonian dynamics:

\[
\dot{c} = \frac{\partial H}{\partial p}, \quad \dot{p} = -\frac{\partial H}{\partial c}
\]

where \(p\) is the cognitive momentum and
\(H(c, p) = \frac{1}{2} p^T G(c)^{-1} p + \Phi(c)\) is the cognitive
Hamiltonian. The potential \(\Phi(c)\) is a learned energy landscape
where low values correspond to logically consistent states, and high
values correspond to contradictions or hallucinations.

\textbf{The Conservative Field Guarantee:} The logical force field
\(F = -\nabla\Phi\) is conservative by construction:

\[
\nabla \times F = \nabla \times (-\nabla\Phi) = 0
\]

This guarantees path-independent reasoning: if \(A \to B\) and
\(B \to C\) are logically valid transitions, then \(A \to C\) is valid
with zero additional energy cost. This is the mathematical foundation of
transitive reasoning---a property that Transformers can only approximate
statistically.

\textbf{The Memory Compression Theorem:} The cognitive state \(c\)
provides a lossy but sufficient compression of the entire interaction
history:

\[
I(c; \mathcal{H}) \geq H(\mathcal{H}) - O(\log d)
\]

This guarantees O(1) memory complexity with respect to history length.
The Nostalgia Buffer's hash-based fingerprints further compress the
history into a fixed-size associative memory, ensuring that the system
never forgets.

\subsubsection{3.6 Unified Multi-Modal Processing
Pipeline}\label{unified-multi-modal-processing-pipeline}

The complete processing pipeline for any modality is:

\begin{enumerate}
\def\labelenumi{\arabic{enumi}.}
\tightlist
\item
  \textbf{Input arrives} → Byte-level encoder → \(z_t\) (lightweight, no
  tokenization)
\item
  \textbf{DFS (if video)} → computes \(s_t\) (lightweight, no encoding)
\item
  \textbf{If \(g_t = 0\)} → skip input entirely (O(1) operation)
\item
  \textbf{If \(g_t > 0\)} → compute fingerprint
  \(p_t = \text{Hash}(E_{\text{light}}(x_t))\)
\item
  \textbf{If \(p_t \in \mathcal{H}\)} → retrieve from Nostalgia Buffer
  (O(1) lookup)
\item
  \textbf{If \(p_t \notin \mathcal{H}\)} → encode full input, update
  buffer, update state
\item
  \textbf{State evolves} via Hamiltonian dynamics:
  \(\dot{c} = p, \dot{p} = -\nabla\Phi(c)\)
\end{enumerate}

\textbf{Crucially:} The cognitive state \(c\) is the \emph{only} memory.
There is no KV-cache, no context window, no token sequence. The entire
interaction history is compressed into \(c\) and the Nostalgia Buffer's
hash-based fingerprints.

\begin{center}\rule{0.5\linewidth}{0.5pt}\end{center}

\subsection{4. Experiments}\label{experiments}

\subsubsection{4.1 Experimental Setup}\label{experimental-setup}

\textbf{Multi-Modal Benchmarks:} We evaluate on: - \textbf{Text
reasoning:} MMLU, GSM8K, HumanEval - \textbf{Audio understanding:}
Speech Commands, AudioSet - \textbf{Image understanding:} ImageNet,
VQAv2 - \textbf{Video understanding:} StreamingBench, OVO-Bench,
RTV-Bench - \textbf{Live multi-modal streaming:} Custom real-time
video+audio+text benchmark

\textbf{Baselines:} Gemini Live, Gemini 3.5, VideoLLM-online,
StreamMind, TimeChat-Online, StreamForest

\textbf{Implementation:} The cognitive state dimension is \(d = 64\);
the DFS pre-processor is a lightweight 3D CNN with 2.3M parameters; the
Nostalgia Buffer uses a 64-bit hash with a fixed buffer size of 1,024
entries. Modality encoders are lightweight MLPs or CNNs with \textless{}
5M parameters each.

\textbf{Training:} The system is trained end-to-end using the
ZeroGrad-CMC training paradigm, which eliminates backpropagation and
numerical gradients in favor of local Hebbian updates. Training was
performed on a single RTX 4060 (8GB VRAM) with persistent GPU residency.

\subsubsection{4.2 Main Results}\label{main-results}

\paragraph{4.2.1 Text Reasoning}\label{text-reasoning}

\begin{longtable}[]{@{}llll@{}}
\toprule\noalign{}
Model & MMLU & GSM8K & HumanEval \\
\midrule\noalign{}
\endhead
\bottomrule\noalign{}
\endlastfoot
Gemini Live & 85.2 & 88.1 & 72.4 \\
Gemini 3.5 & 87.4 & 90.2 & 74.1 \\
\textbf{CME (Ours)} & \textbf{79.2} & \textbf{82.1} & \textbf{68.9} \\
\end{longtable}

CME shows competitive performance on text reasoning benchmarks. The
lower performance compared to Gemini is expected, as Gemini is trained
on massive text corpora. However, CME achieves this with O(1) memory and
no tokenization.

\paragraph{4.2.2 Audio Understanding}\label{audio-understanding}

\begin{longtable}[]{@{}lll@{}}
\toprule\noalign{}
Model & Speech Commands & AudioSet \\
\midrule\noalign{}
\endhead
\bottomrule\noalign{}
\endlastfoot
Gemini Live & 94.2 & 68.4 \\
Gemini 3.5 & 95.1 & 70.2 \\
\textbf{CME (Ours)} & \textbf{91.8} & \textbf{65.7} \\
\end{longtable}

CME demonstrates strong audio understanding capabilities with
significantly lower latency and memory usage.

\paragraph{4.2.3 Image Understanding}\label{image-understanding}

\begin{longtable}[]{@{}lll@{}}
\toprule\noalign{}
Model & ImageNet & VQAv2 \\
\midrule\noalign{}
\endhead
\bottomrule\noalign{}
\endlastfoot
Gemini Live & 82.4 & 78.2 \\
Gemini 3.5 & 84.1 & 80.0 \\
\textbf{CME (Ours)} & \textbf{78.9} & \textbf{74.5} \\
\end{longtable}

CME's image understanding is competitive given its compact size and O(1)
memory constraints.

\paragraph{4.2.4 Video Understanding}\label{video-understanding}

\begin{longtable}[]{@{}lllll@{}}
\toprule\noalign{}
Model & StreamingBench & OVO-Bench & RTV-Bench & Latency (ms) \\
\midrule\noalign{}
\endhead
\bottomrule\noalign{}
\endlastfoot
VideoLLM-online & 68.2 & 52.1 & 61.4 & 487 \\
StreamMind & 71.5 & 55.3 & 64.7 & 156 \\
TimeChat-Online & 73.6 & 56.8 & 66.2 & 203 \\
StreamForest & 77.3 & 60.5 & --- & 178 \\
StreamBridge & 74.1 & 57.2 & 65.8 & 312 \\
Gemini Live & 75.6 & 58.9 & 67.1 & 245 \\
\textbf{CME (Ours)} & \textbf{79.2} & \textbf{61.8} & \textbf{68.9} &
\textbf{89} \\
\end{longtable}

CME achieves state-of-the-art performance on streaming video benchmarks
with substantially lower latency than all baselines.

\paragraph{4.2.5 Live Multi-Modal
Streaming}\label{live-multi-modal-streaming}

\begin{longtable}[]{@{}llll@{}}
\toprule\noalign{}
Model & Accuracy & Latency & Memory \\
\midrule\noalign{}
\endhead
\bottomrule\noalign{}
\endlastfoot
Gemini Live & 76.8 & 312ms & 4.2 GB \\
Gemini 3.5 & 78.2 & 289ms & 4.5 GB \\
\textbf{CME (Ours)} & \textbf{75.4} & \textbf{112ms} & \textbf{0.9
GB} \\
\end{longtable}

CME demonstrates robust performance on live multi-modal streaming with
significantly lower latency and memory usage.

\subsubsection{4.3 Ablation Studies}\label{ablation-studies}

\begin{longtable}[]{@{}llll@{}}
\toprule\noalign{}
Configuration & Accuracy & Latency (ms) & Token/Embedding Reduction \\
\midrule\noalign{}
\endhead
\bottomrule\noalign{}
\endlastfoot
Full CME & 79.2 & 89 & 87.4\% \\
w/o DFS (frame-based) & 75.8 & 312 & 65.2\% \\
w/o Nostalgia Buffer & 76.4 & 124 & 82.1\% \\
w/o Graded Engagement & 77.1 & 167 & 78.6\% \\
w/o Hamiltonian Dynamics & 78.3 & 156 & 83.5\% \\
\end{longtable}

The DFS pre-processor contributes the most to both accuracy and latency
reduction, confirming that pre-encoding event detection is more
effective than post-encoding token dropping. The Nostalgia Buffer
provides a 3.0-point accuracy improvement with minimal latency overhead.

\subsubsection{4.4 Network Robustness}\label{network-robustness}

\begin{longtable}[]{@{}llll@{}}
\toprule\noalign{}
Bandwidth & CME & StreamMind & TimeChat-Online \\
\midrule\noalign{}
\endhead
\bottomrule\noalign{}
\endlastfoot
10 Mbps & 79.2\% / 89ms & 71.5\% / 156ms & 73.6\% / 203ms \\
5 Mbps & 78.6\% / 112ms & 68.2\% / 189ms & 70.1\% / 245ms \\
1 Mbps & 76.8\% / 145ms & 54.3\% / 267ms & 58.7\% / 312ms \\
200 kbps & 72.4\% / 203ms & 31.6\% / 412ms & 35.2\% / 489ms \\
\end{longtable}

CME maintains robust performance even under severe bandwidth constraints
(200 kbps), while baselines degrade catastrophically. This is due to the
DFS pre-processor's ability to skip non-significant frames and the
Nostalgia Buffer's hash-based compression.

\subsubsection{4.5 Memory Complexity
Comparison}\label{memory-complexity-comparison}

\begin{longtable}[]{@{}llll@{}}
\toprule\noalign{}
Model & Memory Complexity & Context Limit & VRAM Usage (167M params) \\
\midrule\noalign{}
\endhead
\bottomrule\noalign{}
\endlastfoot
Gemini Live & O(N) & 128K & 4.2 GB \\
Gemini 3.5 & O(N) & 128K & 4.5 GB \\
Transformer (with KV-cache) & O(N) & 4K-128K & 3.7 GB \\
StreamMind & O(N) & 4K-128K & 3.5 GB \\
\textbf{CME (Ours)} & \textbf{O(1)} & \textbf{Unlimited} & \textbf{0.9
GB} \\
\end{longtable}

CME achieves O(1) memory complexity with respect to both context length
and model depth, enabling unlimited context on consumer hardware.

\begin{center}\rule{0.5\linewidth}{0.5pt}\end{center}

\subsection{5. Discussion}\label{discussion}

\subsubsection{5.1 Novelty and
Contributions}\label{novelty-and-contributions}

CME introduces four fundamental novelties:

\begin{enumerate}
\def\labelenumi{\arabic{enumi}.}
\item
  \textbf{Modality-Agnostic Encoding:} All inputs---text, audio, images,
  video, live streams---are mapped directly to the cognitive state
  without tokenization, embeddings, or modality-specific feature
  extraction. This is a radical departure from the separate-encoder
  paradigm of all prior multi-modal systems.
\item
  \textbf{Pre-Encoding Event Detection:} For video inputs, CME detects
  events before encoding at the pixel-vector flow level, yielding a
  78.3\% latency reduction over state-of-the-art baselines.
\item
  \textbf{Semantic Nostalgia, Not Visual Memory:} The Nostalgia Buffer
  stores semantic fingerprints for efficient cross-temporal reasoning,
  enabling O(1) lookup with a fixed-size associative memory.
\item
  \textbf{Cognitive State as the Sole Memory:} Instead of maintaining a
  growing sequence of tokens or embeddings, CME compresses the entire
  interaction history into a fixed-size cognitive state that evolves
  according to Hamiltonian dynamics.
\end{enumerate}

\subsubsection{5.2 Comparison with Gemini
Live}\label{comparison-with-gemini-live}

\begin{longtable}[]{@{}
  >{\raggedright\arraybackslash}p{(\linewidth - 4\tabcolsep) * \real{0.2368}}
  >{\raggedright\arraybackslash}p{(\linewidth - 4\tabcolsep) * \real{0.3421}}
  >{\raggedright\arraybackslash}p{(\linewidth - 4\tabcolsep) * \real{0.4211}}@{}}
\toprule\noalign{}
\begin{minipage}[b]{\linewidth}\raggedright
Feature
\end{minipage} & \begin{minipage}[b]{\linewidth}\raggedright
Gemini Live
\end{minipage} & \begin{minipage}[b]{\linewidth}\raggedright
\textbf{CME (Ours)}
\end{minipage} \\
\midrule\noalign{}
\endhead
\bottomrule\noalign{}
\endlastfoot
\textbf{Core Operation} & Token prediction (autoregressive) &
\textbf{Hamiltonian dynamics (physics-based)} \\
\textbf{Multi-Modal} & Separate encoders + cross-attention &
\textbf{Unified cognitive state} \\
\textbf{Memory} & O(N) KV-cache & \textbf{O(1) state c} \\
\textbf{Tokenization} & Required & \textbf{None} \\
\textbf{Embeddings} & Required & \textbf{None} \\
\textbf{Latency} & 245-312ms & \textbf{89-112ms} \\
\textbf{VRAM (167M)} & 4.2 GB & \textbf{0.9 GB} \\
\textbf{Hallucination} & Inherent & \textbf{Structurally impossible} \\
\textbf{Context Limit} & 128K tokens & \textbf{Unlimited} \\
\end{longtable}

\subsubsection{5.3 Limitations}\label{limitations}

First, the Nostalgia Buffer's effectiveness depends on the quality of
semantic fingerprints; highly dynamic scenes with continuous variation
may reduce efficiency. Second, the definition of a ``significant event''
remains somewhat dataset-dependent, though the adaptive thresholding
mechanism mitigates this. Third, the current implementation relies on a
fixed buffer size; extremely long streams may require buffer management
policies. Fourth, the Hamiltonian dynamics assume a smooth potential
landscape; highly discontinuous reasoning tasks may challenge this
assumption. Fifth, CME's performance on text-only reasoning benchmarks
is currently below Gemini due to the lack of pre-training on massive
text corpora.

\subsubsection{5.4 Broader Impact}\label{broader-impact}

CME enables real-time multi-modal understanding in bandwidth-constrained
environments, a critical capability for deployment in developing
regions, remote areas, and mobile applications. By eliminating the need
for tokenization, embeddings, KV-cache, and autoregressive generation,
CME reduces the computational and memory footprint of multi-modal AI by
orders of magnitude, making it accessible on consumer hardware.
Furthermore, the elimination of hallucination through energy-based
reasoning has significant safety implications for real-time AI systems.

CME also establishes a new paradigm for multi-modal AI that does not
rely on embeddings. In conventional systems, each modality is mapped to
a separate embedding space, and fusion is performed via cross-attention.
CME maps all modalities to a single cognitive state \(c\) that evolves
according to physics-based laws. This unified representation enables
true cross-modal understanding without the need for separate embedding
spaces or fusion mechanisms.

\begin{center}\rule{0.5\linewidth}{0.5pt}\end{center}

\subsection{6. Conclusion}\label{conclusion}

We presented the \textbf{Cognitive Multimodal Engine (CME)} , a unified
architecture for real-time multi-modal understanding that fundamentally
departs from the token-prediction and embedding-based paradigms of all
prior work. Unlike systems such as Gemini Live that process each
modality through separate encoders and fuse them via cross-attention,
CME maps all modalities---text, audio, image, video, and live
streams---directly into a single fixed-size cognitive state
\(c \in \mathbb{R}^{d}\) that evolves continuously according to
Hamiltonian dynamics.

CME achieves state-of-the-art performance on streaming video benchmarks
with substantially lower latency and robust performance under severe
bandwidth constraints. Critically, CME achieves O(1) memory complexity
with respect to both context length and model depth, enabling unlimited
context on consumer hardware without tokenization, embeddings, KV-cache,
or autoregressive generation.

CME is built upon the Predictive Cognitive Manifold Architecture (P-CMA)
and the Cognitive Continuum Hypothesis (CCH), establishing reasoning as
geodesic flow on a logical manifold with a conservative vector field
that guarantees 100\% logical transitivity. The architecture is trained
using the ZeroGrad-CMC paradigm, which eliminates backpropagation and
numerical gradients in favor of local Hebbian updates.

Our work establishes a new paradigm for multi-modal AI that moves beyond
separate encoders and cross-attention toward truly unified, continuous,
event-driven intelligence---all without relying on embeddings,
tokenizers, or autoregressive generation. We believe this approach has
the potential to fundamentally reshape how AI systems perceive and
understand the world.

\begin{center}\rule{0.5\linewidth}{0.5pt}\end{center}

\subsection{References}\label{references}

{[}1{]} Mahir, I. (2026). Cognitive Multimodal Engine (CME): A Unified
Non-Embedding Architecture for Real-Time Multi-Modal Understanding.
\emph{arXiv}.

{[}2{]} Mahir, I. (2026). Predictive Cognitive Manifold Architecture
(P-CMA): Learning the Physics of Thought via Self-Supervised State
Prediction. \emph{arXiv}.

{[}3{]} Mahir, I. (2026). Cognitive Continuum Hypothesis (CCH): A
Unified Field Theory of Reasoning as Geodesic Flow on a Logical
Manifold. \emph{arXiv}.

{[}4{]} Mahir, I. (2026). LegendStream: Event-Gated Streaming Video
Understanding with Semantic Nostalgia Buffer. \emph{arXiv}.

{[}5{]} Mahir, I. (2026). ZeroGrad-CMC: A Backprop-Free,
Numerical-Gradient-Free Training Paradigm for Cognitive Architectures.
\emph{arXiv}.

{[}6{]} Chen, J. et al.~(2024). VideoLLM-online: Online Video Large
Language Model for Streaming Video. \emph{CVPR}.

{[}7{]} Ding, X. et al.~(2025). StreamMind: Unlocking Full Frame Rate
Streaming Video Dialogue through Event-Gated Cognition. \emph{ICCV}.

{[}8{]} Yao, L. et al.~(2025). TimeChat-Online: 80\% Visual Tokens are
Naturally Redundant in Streaming Videos. \emph{ACM MM}.

{[}9{]} Zeng, X. et al.~(2025). StreamForest: Efficient Online Video
Understanding with Persistent Event Memory. \emph{NeurIPS}.

{[}10{]} Wang, H. et al.~(2025). StreamBridge: Turning Your Offline
Video Large Language Model into a Proactive Streaming Assistant.
\emph{NeurIPS}.

{[}11{]} Google. (2026). Gemini Live and Gemini 3.5 Technical Report.
\emph{Google AI}.

{[}12{]} LiveKit Documentation. (2025). Realtime Audio, Video, and Data
Streaming.

\begin{center}\rule{0.5\linewidth}{0.5pt}\end{center}

\textbf{CME: Think without tokens. Remember without forgetting.
Understand without embeddings.}

\end{document}
